\section{Introduction}\label{sec:introduction}

Let $X$ be a smooth connected complex algebraic curve and let $r\geq1$.
Its $\SL_r$-character variety parametrizes semisimple representations of
$\pi_1(X)$ up to isomorphism. Although $X$ need not be defined over a number
field, this character variety has a natural integral model: a finite
presentation of its fundamental group gives a representation scheme over
$\Z$, and we take the spectrum of its ring of conjugation invariants.
Write this affine scheme as $M_r(X)$.

An arithmetic question is whether its integral points become Zariski dense
after extending the ground field once. Here integrality always refers to
the full ring of integers. We prove this for curves, retaining arbitrary
quasi-unipotent conjugacy classes at the punctures.

Let $\overline X$ be the smooth projective completion of $X$, and write
$\overline X\setminus X=\{x_1,\ldots,x_s\}$. For each $i$, fix a positive
meridian $\gamma_i$ about $x_i$. A matrix is \emph{quasi-unipotent} if all
its eigenvalues are roots of unity. Choose quasi-unipotent conjugacy
classes $C_i\subset\SL_r(\C)$, and let $Y$ be the reduced locus in
$M_r(X)_\C$ whose semisimple representations have
$\rho(\gamma_i)\in C_i$. This condition allows nonsemisimple boundary
matrices: semisimplicity of a representation is a condition on the whole
group action. The locus $Y$ is locally closed, as explained in
Section~\ref{sec:model}.

For a number field $L\subset\C$, write $\cO_L$ for its full ring
of integers.  We use \emph{ambient} integrality on $Y$: a point
extends to $M_r(X)$ over $\cO_L$, and its complex semisimple
representation has the specified boundary classes.  It need not
avoid smaller Jordan classes after reduction at every finite prime.
Requiring a section of a strict locally closed integral model is
a different condition.

\begin{theorem}\label{thm:main}
Let $X$ be a smooth connected complex algebraic curve and let $r\geq1$.
For prescribed quasi-unipotent conjugacy classes $C_1,\ldots,C_s$ as above,
let $K\subset\C$ be a number field containing their eigenvalues. There is
a finite extension $L/K$ such that
\[
 M_r(X)(\cO_L)\cap Y(\C)
\]
is Zariski dense in $Y$, and hence in every irreducible component of $Y$.
The same conclusion holds when only the characteristic polynomial of
each boundary monodromy is prescribed, with all its roots roots of unity.
It also holds for the entire character variety without boundary
conditions; in that case one may take $K=\Q$.

For the relative statements and for projective $X$, one may take distinct
primes $p,q>r$ and
\begin{equation}\label{eq:main-field}
 F=K(\sqrt2,\zeta_{pq}),\qquad
 L=F\bigl(\zeta_r,\{u^{1/r}:u\in\cO_F^\times\}\bigr),
\end{equation}
where $\zeta_m$ denotes a primitive $m$th root of unity. In the projective
case there are no boundary conditions and $K=\Q$.
\end{theorem}

The extension in \eqref{eq:main-field} is finite because the unit group
of $\cO_F$ is finitely generated. It is fixed for the entire locus,
including all its components. No restriction on the genus or on the
Jordan forms is imposed.

For a nonprojective curve without boundary conditions, density already
holds over $\Z$: the fundamental group is free and $\SL_r(\Z)$ is
Zariski dense in $\SL_r(\C)$. The substantive cases here are prescribed
boundary monodromy and projective curves, where the surface relation
couples the matrices. We construct actual integral representations and
then pass to their characters; the exact boundary conditions are selected
only after this passage.

\subsection{History and relation to previous work}

The arithmetic study of local systems includes Simpson's conjecture
that rigid irreducible representations of smooth projective fundamental
groups are defined over rings of integers~\cite[p.~9]{Simpson1992}.
Esnault and Groechenig prove integrality for irreducible cohomologically
rigid local systems with finite determinant and quasi-unipotent boundary
monodromy on smooth quasiprojective varieties~\cite[Theorem~1.1]{EsnaultGroechenig2018}.
Here an integral local system may have a finite projective lattice over a
number ring; freeness is not part of that assertion.
Without a rigidity assumption, de Jong and Esnault prove that the
existence of an irreducible complex local system with torsion determinant
and quasi-unipotent boundary monodromy implies the existence of an
absolutely irreducible integral $\ell$-adic local system for every prime
$\ell$, preserving the determinant and the semisimplified boundary
classes~\cite[Theorem~1.1]{DeJongEsnault2024}.
This prime-by-prime existence is distinct from density over the full ring
of integers of one number field.

Litt asks for potential integral density on character varieties of smooth
projective varieties and also formulates a boundary
variant~\cite{LittProblems}; his survey explicitly asks whether one
number field suffices~\cite[Question~5.4.3(2)]{Litt2024}.
Coccia and Litt formulate a related
conjecture for Chevalley groups with fixed boundary data in the adjoint
quotient~\cite[Conjecture~1.1.1]{CocciaLitt}. Theorem~\ref{thm:main}
answers the determinant-one form of Litt's question affirmatively for
curves, in every rank, and proves the $\SL_r$ curve case of the
Coccia--Litt conjecture for quasi-unipotent boundary data.
The higher-dimensional, higher-rank question remains outside its scope.

Coccia and Litt prove full-ring potential integral density for $\SL_2$
and $\mathrm{PGL}_2$ on smooth varieties admitting a smooth projective
compactification with simple normal crossings boundary~\cite[Theorem~1.1.2]{CocciaLitt}.
Their surface theorem allows arbitrary algebraic-integer boundary traces
and gives an extension of degree at most four over the field generated
by those traces~\cite[Theorem~5.0.4 and Remark~5.0.5]{CocciaLitt}.
Their ordinary and twisted surface theorems also imply the ambient-integral
exact-class statement in rank two, as explained in
Remark~\ref{rem:rank-two-comparison}.
Their use of integral representations and mapping-class-group dynamics
is an arithmetic predecessor of the closed-surface argument below.
Theorem~\ref{thm:main} treats every rank with exact quasi-unipotent
boundary data; in rank two, Coccia and Litt allow a broader range of
boundary traces.

Earlier arithmetic work includes Whang's descent framework for
rank-two integral characters with fixed integer boundary
traces~\cite{WhangDescent}. For the Markoff cubic family
$x^2+y^2+z^2-xyz=k$, Ghosh and Sarnak prove integral density for
almost all locally admissible integer parameters, and also exhibit
infinitely many admissible parameters with no integral
points~\cite[Theorem~1.2]{GhoshSarnak}. These results concern ordinary
integral points and help distinguish that question from density
after one finite extension of the ground field.

On the geometric side, Mellit uses flags and braid operations to stratify
a vector bundle over a character variety into torus and affine-space
pieces, for generic semisimple data with a regular semisimple
puncture~\cite[Theorem~1.2]{Mellit}. Su constructs an open dense torus
for nonempty character varieties with generic semisimple boundary data
and a regular semisimple puncture~\cite[Theorem~0.5]{Su}.
These results concern complex geometric decompositions. The present
argument must also preserve integral splittings and identified graded
modules while treating arbitrary Jordan types and nongeneric components.
It therefore carries out the local operations over the full ring of
integers, rather than obtaining integral points merely from a complex
coordinate parametrization.

\subsection{The proof}

We first work with $\GL_n$ rather than $\SL_n$. This makes direct sums,
subspaces and quotients available without determinant restrictions.
A boundary matrix with a prescribed Jordan type admits an invariant
flag on whose successive quotients it acts by scalar roots of unity.
We attach disks carrying these scalar local systems to the boundary.
The resulting object consists of oriented surfaces of various ranks,
joined along boundary pieces by flags and their graded identifications.
Section~\ref{sec:diagrams} defines exactly which joins are allowed.

The main construction proves integral density for all such diagrams by
induction on the largest rank. A surface of maximal rank with boundary
is cut into disks. Two flags at the ends of a cut are compared by
refining them to their common graded pieces and interchanging adjacent
pieces. On each disk, a proper subspace chosen from one boundary flag
is preserved by all puncture monodromies, since these are scalar. Passing
to that subspace and its quotient lowers the rank. The lost extension
data are recovered by successive choices of linear maps. Choosing the
subspace as a prefix of the starting flag makes the lifted flag unique
on that interval, so returning around the boundary introduces no further
equation.
The reduction includes the complementary swaps needed for
nonregular boundary data (Section~\ref{sec:reduction}).

Closed surfaces enter at the same induction step. Cutting along one
nonseparating circle and fixing a convenient spectrum brings them into
the preceding case. Powers of a Dehn twist then enlarge the closure of
the integral points. At one explicit smooth representation, the
resulting torus motions vary the spectrum in every direction.
Irreducibility of the twisted surface equation finishes the argument.
We give a finite-field proof of irreducibility, following the
character-counting method of Liebeck and Shalev~\cite{LiebeckShalev2005}.
Sections~\ref{sec:surfaces} and~\ref{sec:dynamics} carry out this
closed-surface step; the latter then assembles the rank induction.

Finally, all handle determinants of our integral representations are
units in $\cO_F$. Adjoining their $r$th roots uniformly by
\eqref{eq:main-field} permits determinant-one twisting. Exact Jordan
conditions are then selected on the character quotient. This order is
essential: taking the semisimplification of a representation can reduce
its boundary Jordan blocks. Section~\ref{sec:descent} completes this descent.

The reusable ingredient is the diagram density theorem. It transfers
integral density through rank reduction using free graded modules and
linear extension parameters.

Three standard inputs enter at distinct points.  Finite generation
of integral invariants gives the affine model in
Section~\ref{sec:model}.  Dirichlet's unit theorem gives the single
field in Section~\ref{sec:arithmetic}.  The Lang--Weil estimate
converts the finite-group count proved in Section~\ref{sec:surfaces}
into geometric irreducibility.  The flag reconstruction, character
estimates, twist directions and exact-Jordan descent are proved below.
